% !TeX root = ../../../Book1.tex
\section{Lebesgue 可测集}
\label{b1:sec:0302}

将 \(m^*\) 代入 Carathéodory 判别，便得到欧氏空间上最常用的完备测度。它包含全部 Borel 集，也额外包含 Borel 零集的任意子集。

\subsection{Carathéodory 可测性}
\label{b1:subsec:030201}

\begin{definition}
\label{b1:def:lebesgue-measurable}
集合 \(E\subseteq\mathbb R^d\) 称为\term[lebesgue-kekong]{Lebesgue 可测的}[Lebesgue measurable]，若对任意 \(A\subseteq\mathbb R^d\)，
\[
m^*(A)=m^*(A\cap E)+m^*(A\setminus E).
\]
全体 Lebesgue 可测集组成的 \(\sigma\)-代数记为 \(\mathcal L^d\)，而
\[
m=m^*|_{\mathcal L^d}
\]
称为 \(\mathbb R^d\) 上的 Lebesgue 测度。
\end{definition}

由 \cref{b1:thm:caratheodory}，\((\mathbb R^d,\mathcal L^d,m)\) 是完备测度空间。这里的完备性不是附加约定，而是外测度构造自动给出的结果。

\subsection{Borel 集的可测性}
\label{b1:subsec:030202}

\begin{theorem}
\label{b1:thm:borel-lebesgue}
\[
\mathcal B(\mathbb R^d)\subseteq\mathcal L^d.
\]
特别地，所有开集、闭集、半开长方体和连续实函数的反像集都是 Lebesgue 可测集。
\end{theorem}
\begin{proof}
由 \cref{b1:prop:rectangle-outer-measure} 及第二章的 Carathéodory 延拓定理，半开长方体都是 \(m^*\)-可测集。它们生成 \(\mathcal B(\mathbb R^d)\)，而 \(\mathcal L^d\) 是 \(\sigma\)-代数，故包含所有 Borel 集。连续函数的开集反像是开集，故其 Borel 原像也 Lebesgue 可测。
\end{proof}

\subsection{零测集}
\label{b1:subsec:030203}

\begin{definition}
\label{b1:def:lebesgue-null}
若 \(E\subseteq\mathbb R^d\) 满足 \(m^*(E)=0\)，则称 \(E\) 为\term[lingceji]{Lebesgue 零测集}[Lebesgue null set]。
\end{definition}

\begin{proposition}
\label{b1:prop:null-subsets}
若 \(N\) 是 Lebesgue 零测集，且 \(E\subseteq N\)，则 \(E\in\mathcal L^d\) 且 \(m(E)=0\)。
\end{proposition}
\begin{proof}
由单调性 \(m^*(E)\leq m^*(N)=0\)。\cref{b1:rem:outer-measure-null} 表明每个外测度为零的集合都 Carathéodory 可测，故结论成立。
\end{proof}

例如任意单点在 \(\mathbb R^d\) 中零测：用边长趋于零的立方体覆盖它即可。可数集因而也零测。

\subsection{Lebesgue \texorpdfstring{\(\sigma\)}{σ}-代数}
\label{b1:subsec:030204}

\begin{theorem}
\label{b1:thm:lebesgue-completion-borel}
\(\mathcal L^d\) 是 Borel \(\sigma\)-代数关于 Lebesgue 测度限制的完备化。等价地，
\[
\mathcal L^d
=\{B\cup N:B\in\mathcal B(\mathbb R^d),\
N\subseteq Z\text{ 且 }Z\in\mathcal B(\mathbb R^d),\ m(Z)=0\}.
\]
\end{theorem}
\begin{proofsketch}
Borel 集由半开长方体生成，故已包含于 \(\mathcal L^d\)；\cref{b1:prop:null-subsets} 再补入全部 Borel 零集的子集。反向包含可通过外测度覆盖构造：对 \(E\in\mathcal L^d\)，选开集 \(G_n\supseteq E\)，使 \(m(G_n\setminus E)<2^{-n}\)，并取适当的 Borel 交、并组合，得到一个 Borel 集 \(B\) 与 \(E\) 的对称差零测。
\end{proofsketch}
